\documentclass[12pt, a4paper]{article}

% --- Professional Scientific Preamble ---
\usepackage[utf8]{inputenc}
\usepackage[T1]{fontenc}
\usepackage{amsmath, amssymb, amsfonts} 
\usepackage{booktabs}                   
\usepackage{geometry}                   
\usepackage{url}                        
\usepackage{enumitem}                   
\geometry{margin=1in}

% --- Document Metadata ---
\title{Mathematical Topology and Harmonic Approximation of 72-EDO: \\ \large A Study in High-Resolution Temperament}
\author{Acoustical Research Group}
\date{\today}

\begin{document}

\maketitle

\begin{abstract}
This paper provides a formal analysis of 72-EDO (72 Equal Divisions of the Octave). We evaluate its capacity for 11-limit Just Intonation (JI) approximation and its topological efficiency using Fokker Periodicity Blocks. Our findings confirm 72-EDO as an optimal superset for unifying Western chromaticism and prime-limit microtonality.
\end{abstract}

\section{Logarithmic Foundation}
The system is defined by a step size $\Delta$ of exactly $16.\overline{6}$ cents, derived from the $72^{nd}$ root of 2:
\begin{equation}
f(n) = f_0 \cdot 2^{n/72}, \quad n \in \{0, 1, \dots, 71\}
\end{equation}

\section{Harmonic Mapping and Error Metrics}
72-EDO is mathematically superior to 12-EDO because it captures the $5^{th}, 7^{th},$ and $11^{th}$ harmonics with high fidelity. The "truth" of the system is demonstrated in its minimal deviation ($\epsilon$) from the physical harmonic series.



\begin{table}[h!]
\centering
\caption{72-EDO Prime-Limit Mapping}
\label{tab:harmonics}
\begin{tabular}{@{}llcccr@{}}
\toprule
Harmonic & Ratio & Just (Cents) & 72-EDO Step & 72-EDO (Cents) & Error (Cents) \\ \midrule
3  & $3/2$   & 701.96  & 42 & 700.00  & $-1.96$ \\
5  & $5/4$   & 386.31  & 23 & 383.33  & $-2.98$ \\
7  & $7/4$   & 968.83  & 58 & 966.67  & $-2.16$ \\
11 & $11/8$  & 551.32  & 33 & 550.00  & $-1.32$ \\
13 & $13/8$  & 840.53  & 50 & 833.33  & $-7.20$ \\ \bottomrule
\end{tabular}
\end{table}

\section{Fokker Periodicity Blocks and Commas}
The structural "truth" of 72-EDO lies in its comma resolution. Unlike 12-EDO, which tempers out the Syntonic Comma ($81/80$), 72-EDO maps it to exactly 1 step (approx. $16.67$ cents). This allows for the construction of Fokker Periodicity Blocks \cite{erlich1998}, which represent JI intervals as a closed lattice in a finite equal temperament.



\section{Historical and Analytical References}
The following sources provide the foundational "truth" for this study.

\begin{thebibliography}{99}

\bibitem{monzo2005}
Joe Monzo. \textit{72-EDO}. Tonalsoft Encyclopedia of Microtonal Music Theory, 2005. \\
Available: \url{http://www.tonalsoft.com/enc/number/72edo.aspx}

\bibitem{erlich1998}
Paul Erlich. \textit{A Gentle Introduction to Fokker Periodicity Blocks}. Xenharmonikon 17, 1998. \\
(Crucial for understanding the geometric mapping of 72-EDO).

\bibitem{wysch1937}
Ivan Wyschnegradsky. \textit{Manuel d'harmonie à quarts de ton}. Paris, 1937. \\
(Historical foundation of 72-EDO as a superset of smaller EDOs).

\bibitem{maneri1994}
Joe Maneri and Scott Van Duyne. \textit{Preliminary Studies in the Virtual Pitch Continuum}. Boston, 1994. \\
(Reference for 72-EDO notation and pedagogical application).

\bibitem{partch1974}
Harry Partch. \textit{Genesis of a Music}. Da Capo Press, 1974. \\
(Foundational text for the 11-limit JI theory that 72-EDO approximates).

\end{thebibliography}

\end{document}
